\section{Research agenda: precise next theorems and experiments}
\label{sec:research}

The next stage should prioritize questions that either remove an identified
asymptotic obstruction or determine whether the new kernels accelerate
actual difficult inputs. The following questions are deliberately stated
with success criteria. A negative answer can also be valuable: an explicit
obstruction prevents effort being spent on an insufficient invariant or
representation. These are proposed research problems, not claims that the
necessary bounds have been proved.

\subsection{Near-term algebra and implementation}

\begin{researchquestion}[Realization of directional separation]
Can a family of actual knot diagrams and legal scan prefixes realize the
shared-suffix transfer structure of Theorem~\ref{thm:corridor-separation},
up to uniformly bounded local modifications?
\end{researchquestion}

The algebraic construction has fixed matching width and a nonzero
three-dot output, so its separation is not an artifact of an invalid
differential or zero answer. Its remaining weakness is geometric
realization. A useful positive result would give explicit diagrams,
crossing orders, and a proof identifying the relevant scalar contraction
with the constructed family. One should measure the complete scan as well
as the isolated stage, and compare with the production minimum-fill
algorithm. A useful negative result would characterize a restriction on
scalar contractions of actual prefixes that rules this family out. Random
diagram searches cannot establish either asymptotic statement.

\begin{researchquestion}[A proved adaptive switching rule]
Can one choose between sparse Gaussian elimination, forward transfer, and
reverse transfer with a provable bound relative to the best of the three,
including the cost of estimating and constructing the candidate reductions?
\end{researchquestion}

The current adaptive budget permits some sparse work and then builds a
complete corridor. The six ordinary fixtures did not trigger a switch.
Thus that experiment supplies no evidence that the threshold identifies
difficult stages. A promising implementation experiment is to maintain
cheap online quantities already exposed by scalar pivots: predicted Schur
update count, scalar component sizes, source/target imbalance, and observed
coefficient support. The target is a bounded-regret accounting theorem for
the stage, or a carefully charged portfolio implementation. Paying for two
large representations before making the decision would defeat this goal.

\begin{researchquestion}[Scalar bases that minimize transfer support]
What can be proved about minimizing
\[
  |\operatorname{supp}i|+|\operatorname{supp}p|
  +|\operatorname{supp}h|+\sum_j\min(F_j,B_j)
\]
over typed special contractions of a fixed scalar complex?
\end{researchquestion}

The graded survivor multiplicities are fixed, but all four displayed
support quantities depend on the basis. This separates an invariant
obstruction from a representation choice. First classify tractable scalar
graphs, such as forests, bounded-width decompositions, or components of
bounded size; then compare canonical local bases with minimum-fill choices.
A hardness proof for the unrestricted optimization would also be useful,
provided it respects the chain condition and the permitted changes of
matching and shifts. The delivered component contraction is an exact
baseline because it preserves the predecessor's basis; it is not proposed
as a globally optimal basis.

\begin{researchquestion}[Typed support pruning with a favorable cost bound]
Can the homogeneous support condition in
Proposition~\ref{prop:corridor-typed} be maintained during graph evaluation
at less cost than the coefficient work it eliminates on a meaningful input
class?
\end{researchquestion}

This condition is proved but is not the pruning policy measured in the
delivered corridor experiments. Its implementation would maintain source
and target matching/shift types, rather than blindly applying a pairwise
test for every possible endpoint. Grouping ports by type and shift may
reduce that cost. The existing support audit found zero additional
zero-differential shortcuts beyond the degree-gap certificate, so a new
experiment must look for savings inside a nonzero transfer, not count the
same already certified stages again. Both rejected products and setup
operations should be recorded.

\begin{researchquestion}[Graded transfer inside an exact homological window]
Can the existing guard-degree window construction retain enough quantum
data to support a corridor contraction without introducing dependence on
objects that the window deliberately omitted?
\end{researchquestion}

This is an integration theorem, not just another command-line choice.
One needs a commuting restriction statement: the transferred differential
on each queried degree must agree with that of the full complex, with
precisely specified guard layers. The proof must address scalar homotopy
paths that cross the nominal degree boundary. It should then provide an
allocation bound for the retained graded objects and preserve the current
resource-limit semantics. Until that theorem and its implementation exist,
the delivered API rejects corridor reduction together with window probes.

\subsection{Structural algebraic questions controlling global complexity}

\begin{researchquestion}[A universal multiplicity bound or a succinct substitute]
Does every $n$-crossing knot diagram admit a computable decomposition for
which the necessary scalar survivor data have total size
$\exp(O(\log^2 n))$? If explicit survivors are too numerous, can they be
represented and contracted succinctly within that bound?
\end{researchquestion}

This is the central scanner question. Counting noncrossing matching types
does not count copies of a matching, and a boundary-width bound does not
automatically control the homology of the scalar part. An adequate theorem
must bound intermediate multiplicities or replace their explicit storage
by a representation with proved costs for crossing insertion, composition,
cancellation, and the final decision. Merely compressing equal objects in
a table while retaining all matrix entries leaves the same possible
obstruction. The project's existing component-sharing and continuation
work provide relevant baselines that any proposed representation must
surpass on a specified family.

\begin{researchquestion}[Computing only the information required for recognition]
Can an exact continuation quotient or a dual certificate detect the unknot
without materializing the complete intermediate or final Khovanov complex,
while preserving correctness under every allowable remaining closure?
\end{researchquestion}

Full homology can contain information unnecessary for the Boolean
recognition problem. That observation becomes useful only with an exact
equivalence relation on partial states and a bounded representation of the
equivalence classes. A candidate quotient must retain all distinctions
observable by the remaining tangle; equality of current ranks or Euler
characteristics is insufficient. A constructive target is a finite
presentation of the relevant continuation maps, together with a proof
that their number or descriptive complexity stays quasi-polynomial.
Equally useful would be a family forcing many independent continuation
responses, delimiting which quotients can succeed.

\begin{researchquestion}[Transfer with internal separators]
Can a survivor corridor be evaluated through small internal interfaces,
rather than from every source or every target, with a bound in terms of
interface size and coefficient type complexity?
\end{researchquestion}

The delivered algorithm selects one direction per weak component. A
component can still have many sources and many targets even when most
paths pass through a small interface. One may try to factor its path-sum
matrix through that interface, retaining typed composition order and
shared intermediate expressions. A rigorous theorem must charge matrix
fill, coefficient growth, and the possibly dense required output. The
output-size barrier suggests a factored representation as the natural
goal. Its subsequent gluing and scalar reduction would then need their
own exact algorithms; evaluating every output entry immediately would
discard the possible gain.

\subsection{Geometric extensions and hierarchy search}

\begin{researchquestion}[Recognizing the dihedral presentation from topology]
Given binary normal-surface coordinates or interval-pairing data for a
parallelity region, can one construct, or reject, a compatible global
dihedral covering presentation in polynomial bit complexity, with a
checkable certificate relating it to the original embedded surface?
\end{researchquestion}

The new theorem begins after this presentation has been supplied. A
certificate should identify the smaller base, give its peripheral words,
choose a common cyclic coordinate on the fibre, and verify every gluing
map in that coordinate. Locally affine descriptions do not ensure that
all pieces admit the same global coordinate. Boundaries, omitted intervals,
and changes of local sheet order are part of the problem. An extraction
algorithm should report its failure to meet the promise explicitly, so
that the general compressed topology method remains available.

There is an elementary obstruction to refining every marked interval
until it is invariant under monodromy. For the single map
$x\mapsto x+1\pmod W$, the orbit of any one marked breakpoint has
exactly $W$ distinct positions: equality after $j$ steps first occurs
when $W$ divides $j$. Such a refinement can therefore have exponential
size in $\log W$ even for this simplest dihedral presentation. An
extraction algorithm must keep breakpoint orbits and attachment
distinctions symbolically instead of expanding them.

\begin{researchquestion}[Piecewise dihedral monodromy]
What compact classification survives if each monodromy map is a union of
$p$ affine pieces $x\mapsto\pm x+a$, on binary-encoded intervals, rather
than a single global map on $\mathbb Z/W\mathbb Z$?
\end{researchquestion}

The three-family proof uses one translation subgroup and a single
involution on its residue classes. Piecewise maps need not preserve that
quotient. A productive first problem is to bound the number of distinct
component descriptions in terms of $p$ and the base complexity, or to give
an explicit family disproving such a bound for a proposed format. The
appropriate goal may be a small program answering component and boundary
queries, instead of a constant number of surface records. Any extension
must retain orientation parity and boundary cycle lengths, not merely
the total number of orbits.

\begin{researchquestion}[Attachment-aware equivalence of compressed components]
When can two isomorphic cover components be merged as states of a
three-manifold hierarchy after their attachments to the remaining handles
and boundary pattern are included?
\end{researchquestion}

The exact type-count corollary classifies covers over the base. It does not
identify all embeddings or marked attachments. Even a fixed cover may
admit many inequivalent markings. A useful theorem would describe the
automorphisms relevant to the external attachment data and give a
canonical orbit representative, with polynomial or quasi-polynomial bit
complexity. The supplied marked-sheet queries are a small necessary
interface for this task. They allow named attachments to be located but
do not enumerate or compare arbitrary attachment systems.

\begin{researchquestion}[Bounding hierarchy restarts]
Can the restart analysis be replaced by an amortized potential or a state
equivalence for which the number of distinct necessary states is
$\exp(O(\log^2 n))$, even when the hierarchy depth is linear in $n$?
\end{researchquestion}

The arithmetic in Section~\ref{sec:accounting} identifies the obstruction:
a decreasing mixed-radix integer can still have exponentially many
values. Compressing each state is not a bound on how many states an
algorithm visits. A successful new potential must constrain which
patterns can actually follow one another, or show that many decreases
can be skipped through a certified transition. Another possible route is
to merge equivalent restart states while retaining enough peripheral and
attachment information to justify every subsequent operation. An
example forcing exponentially many distinct states for a proposed
potential would precisely identify what must be added.

\begin{researchquestion}[Constructing hierarchies rather than checking them]
Which search operation in the proposed geometric recognizer can be
implemented within a proved quasi-polynomial budget, under exactly the
irreducibility and boundary-pattern hypotheses available at that step?
\end{researchquestion}

The desired result must include the search for a suitable decomposing
surface or certified alternative, not just the verification of a supplied
hierarchy. A promising development would isolate one search subproblem
and give a full input/output contract, encoding bound, and failure
certificate. The geometric cover kernel can then be a named subroutine
with established bit complexity. This modular approach makes it possible
to identify the first still-unproved operation instead of absorbing it
into an unspecified compressed-topology oracle.

\subsection{Verification as a research instrument}

\begin{researchquestion}[Small independently checkable proof objects]
Can both kernels emit certificates substantially cheaper to verify than
recomputing their outputs, and can their finite algebraic cores be
formalized in the project's intended proof environment?
\end{researchquestion}

For the geometric kernel, natural certificates are gcd identities,
parity consistency equations, and boundary permutation formulas. For
scalar transfer they are local contraction identities and a typed
acyclic factorization of the survivor map. Dense global certificates may
be too large even when the computation is succinct, so verification cost
belongs in the theorem statement. The present tests deliberately use
independent expanded-sheet and coefficient-level oracles. They support
implementation confidence but do not replace a formal proof or validate
all future changes automatically.

\paragraph{Suggested order of work.}
The immediate engineering priorities are the realization search,
adaptive cost model, and graded-window compatibility, because they can
be judged on exact outputs and actual difficult diagrams. In parallel,
the dihedral extraction problem can turn the standalone geometric
theorem into a usable hierarchy primitive. The multiplicity/continuation
questions and restart-state bound deserve the main long-term theoretical
effort: they address the obstructions that local constant-factor
improvements cannot remove. Future reports should identify which of
these obligations a result discharges and retain failed experiments
that materially change the direction of research.
